Reasoning traces are now a standard test-time control knob for large language models, but it is still unclear whether longer traces reliably improve generalization under distribution shift. We run a controlled study that isolates trace policy effects using real \gptfourone API calls, fixed decoding settings, and a shared evaluation harness. We compare four fixed policies (\none, \short, \medium, \longp) and one uncertainty-triggered \adaptive policy on an \id/\ood benchmark mix (\gsm as \id; \arc, \bbh tasks, and \mathfive as \ood).

Across 48 examples (18 \id, 30 \ood), all policies reach 1.00 \id accuracy, so robustness differences are entirely driven by \ood behavior. \adaptive yields the best \ood accuracy (0.800), the smallest robustness gap (0.200), and the best \ood calibration (\ece 0.1823). The no-trace regime collapses to 0.367 \ood accuracy, while the longest fixed regime is expensive and underperforms \medium and \adaptive. Pairwise testing shows \adaptive significantly outperforms \none (+0.433, 95\% bootstrap CI [0.233, 0.633], BH-corrected $q=0.0078$), with positive but non-significant gains over stronger fixed baselines at this sample size.

These results support a non-monotonic length--robustness pattern and suggest that confidence-triggered trace budgeting is a practical default for deployment when robustness, calibration, and token cost must be balanced.
